\begin{lemma}
\label{lemma-constant-sheaf-ext-stalk-spectral-sequences}
\begin{reference}
\cite[Chapter IV, \S 4.3, Remark 2, p. 193]{Tohoku}
\end{reference}
Let $X$ be a topological space, let $R$ be a ring, let $M$ be an
$R$-module, let $\mathcal{A}$ be an $\underline{R}$-module on $X$, and let
$x \in X$. Denote by $\mathcal{N}_x$ the filtered set of open
neighbourhoods of $x$, ordered by reverse inclusion. There is a functorial
first quadrant spectral sequence
\begin{equation}
\label{equation-constant-sheaf-ext-stalk-first-spectral-sequence}
{}^{\mathrm{II}}E_2^{p, q} =
\colim_{U \in \mathcal{N}_x}
\Ext_R^p(M, H^q(U, \mathcal{A}))
\quad\Longrightarrow\quad
\SheafExt^{p + q}_{\underline{R}}(\underline{M}, \mathcal{A})_x.
\end{equation}
It induces a finite filtration on the abutment in every total degree.

Let $I \subset R$ be a two-sided ideal, put $S = R/I$, and assume that
$\mathcal{A}$ is an $\underline{S}$-module, viewed as an
$\underline{R}$-module. There is also a functorial first quadrant spectral
sequence
\begin{equation}
\label{equation-constant-sheaf-ext-stalk-quotient-spectral-sequence}
{}^{\mathrm{III}}E_2^{p, q} =
\bigoplus_{a + b = q}
\colim_{U \in \mathcal{N}_x}
\Ext_S^p(\text{Tor}_a^R(S, M), H^b(U, \mathcal{A}))
\quad\Longrightarrow\quad
\SheafExt^{p + q}_{\underline{R}}(\underline{M}, \mathcal{A})_x.
\end{equation}
It too induces a finite filtration in every total degree. Both spectral
sequences are additive and functorial in $\mathcal{A}$ from the $E_2$-page
onward.
\end{lemma}

\begin{proof}
Choose once and for all a free resolution $L_\bullet \to M$ and the
Godement resolution
$\mathcal{A} \to \mathcal{C}^\bullet(\mathcal{A})$ of
Lemma \ref{lemma-godement-resolution}. The construction of the Godement
resolution commutes with restriction to open subspaces: this follows
directly from its construction by the discrete space having the same
points. Thus, for $U \in \mathcal{N}_x$, the complex
$$
C_U^\bullet = \Gamma(U, \mathcal{C}^\bullet(\mathcal{A}))
$$
computes $R\Gamma(U, \mathcal{A})$, and an inclusion $V \subset U$ gives
a restriction map $C_U^\bullet \to C_V^\bullet$. These maps commute
strictly with composition.

Apply Lemma \ref{lemma-constant-sheaf-ext-double-complex} on every $U$,
using the fixed resolution $L_\bullet$. The restriction maps just described
give a directed system, from the $E_2$-page onward, of first quadrant
spectral sequences with
$$
E_2^{p, q}(U) = \Ext_R^p(M, H^q(U, \mathcal{A}))
$$
and abutment
$$
H_U^n =
\Ext^n_{D(\underline{R}|_U)}
(\underline{M}|_U, \mathcal{A}|_U).
$$
For the quotient construction, apply
Lemma \ref{lemma-constant-sheaf-ext-quotient-spectral-sequence} on every
$U$. Its functoriality with respect to the maps
$C_U^\bullet \to C_V^\bullet$ gives a directed system from the $E_2$-page
onward. Its $E_2$-page is
$$
\bigoplus_{a + b = q}
\Ext_S^p(\text{Tor}_a^R(S, M), H^b(U, \mathcal{A})).
$$
The canonical comparison maps between the Cartan--Eilenberg resolutions
used in that lemma are functorial from the $E_2$-page onward, so the maps
for restrictions are compatible with identities and composition.

We explain carefully the passage to the filtered colimit. For either
directed system write $E_r(U)$ for its $r$th page and set
$$
\overline E_r^{p, q} =
\colim_{U \in \mathcal{N}_x} E_r^{p, q}(U),
\qquad
\overline d_r = \colim_{U \in \mathcal{N}_x} d_r(U).
$$
Filtered colimits of abelian groups are exact. Consequently they commute
with kernels, images, and cokernels, and hence
$$
H(\overline E_r, \overline d_r) = \overline E_{r + 1}.
$$
Thus the colimit pages form spectral sequences with the displayed
$E_2$-pages. In the quotient spectral sequence the sum indexed by
$a + b = q$ is finite for fixed $q$, so it commutes with the filtered
colimit and gives exactly the second displayed $E_2$-page.

For completeness, convergence also passes to the colimit. In total degree
$n$, each spectral sequence over $U$ has a finite filtration
$$
H_U^n = F_U^0H_U^n \supset F_U^1H_U^n \supset \ldots
\supset F_U^{n + 1}H_U^n = 0
$$
whose graded pieces are $E_\infty^{p, n - p}(U)$. Restriction maps preserve
these filtrations. Exactness of filtered colimits gives exact sequences
$$
0 \longrightarrow \colim_U F_U^{p + 1}H_U^n
\longrightarrow \colim_U F_U^pH_U^n
\longrightarrow \colim_U E_\infty^{p, n - p}(U)
\longrightarrow 0.
$$
Hence $F^p = \colim_U F_U^pH_U^n$ is a finite filtration of
$\colim_U H_U^n$. For a fixed bidegree $(p, q)$, a first quadrant spectral
sequence stabilizes after a page depending only on $p$ and $q$, and not on
$U$. Consequently
$\colim_U E_\infty^{p, q}(U) = \overline E_\infty^{p, q}$, so the graded
pieces of this filtration are the stable page of the colimit spectral
sequence.

Finally, by Lemma \ref{lemma-section-RHom-over-U}, the sheaf
$\SheafExt^n_{\underline{R}}(\underline{M}, \mathcal{A})$ is the
sheafification of the presheaf $U \mapsto H_U^n$. Sheafification does not
change stalks, and therefore
\begin{equation}
\label{equation-constant-sheaf-ext-stalk-colimit}
\SheafExt^n_{\underline{R}}(\underline{M}, \mathcal{A})_x
= \colim_{U \in \mathcal{N}_x}
\Ext^n_{D(\underline{R}|_U)}
(\underline{M}|_U, \mathcal{A}|_U).
\end{equation}
This identifies the colimit abutments and their finite filtrations with
those asserted in the lemma. Additivity and functoriality pass to filtered
colimits.
\end{proof}


\begin{lemma}
\label{lemma-constant-sheaf-ext-stalk-comparison}
\begin{reference}
\cite[Chapter IV, \S 4.3, Remark 2, p. 193]{Tohoku}
\end{reference}
In the situation of
Lemma \ref{lemma-constant-sheaf-ext-stalk-spectral-sequences}, there are
canonical maps
\begin{equation}
\label{equation-constant-sheaf-ext-stalk-comparison}
\SheafExt^n_{\underline{R}}(\underline{M}, \mathcal{A})_x
\longrightarrow \Ext_R^n(M, \mathcal{A}_x),
\qquad n \geq 0.
\end{equation}
They are isomorphisms if either one of the following sets of hypotheses
holds.
\begin{enumerate}
\item For every $p, q \geq 0$, the canonical map
$$
\colim_{U \in \mathcal{N}_x}
\Ext_R^p(M, H^q(U, \mathcal{A}))
\longrightarrow
\Ext_R^p\left(M,
\colim_{U \in \mathcal{N}_x}H^q(U, \mathcal{A})\right)
$$
is an isomorphism.
\item There is a two-sided ideal $I \subset R$ such that, with $S = R/I$,
the sheaf $\mathcal{A}$ is an $\underline{S}$-module, and for every
$a, b, p \geq 0$ the canonical map
$$
\begin{aligned}
&
\colim_{U \in \mathcal{N}_x}
\Ext_S^p(\text{Tor}_a^R(S, M), H^b(U, \mathcal{A}))
\\
&\qquad\longrightarrow
\Ext_S^p\left(\text{Tor}_a^R(S, M),
\colim_{U \in \mathcal{N}_x}H^b(U, \mathcal{A})\right)
\end{aligned}
$$
is an isomorphism.
\end{enumerate}
\end{lemma}

\begin{proof}
The exact stalk functor sends $\underline{M}|_U$ to $M$ and
$\mathcal{A}|_U$ to $\mathcal{A}_x$. It therefore induces maps
$$
\Ext^n_{D(\underline{R}|_U)}
(\underline{M}|_U, \mathcal{A}|_U)
\longrightarrow \Ext_R^n(M, \mathcal{A}_x)
$$
compatible with restriction in $U$. Their colimit, under the identification
of Equation (\ref{equation-constant-sheaf-ext-stalk-colimit}), is the map
in Equation (\ref{equation-constant-sheaf-ext-stalk-comparison}).

We first record the cohomology groups appearing after passage to the stalk:
$$
\colim_{U \in \mathcal{N}_x} H^q(U, \mathcal{A}) =
\begin{cases}
\mathcal{A}_x & q = 0, \\
0 & q > 0.
\end{cases}
$$
For $q = 0$ this is the definition of the stalk. If $q > 0$, then every
class over a neighbourhood $U$ vanishes after restriction to an open
covering of $U$ by
Lemma \ref{lemma-kill-cohomology-class-on-covering}. Choosing a member of
that covering which contains $x$ shows that the germ of the class is zero.

The fixed Godement resolution gives strict maps
$C_U^\bullet \to \mathcal{C}^\bullet(\mathcal{A})_x$. By
Lemma \ref{lemma-godement-resolution}, the augmentation
$\mathcal{A}_x[0] \to \mathcal{C}^\bullet(\mathcal{A})_x$ is a homotopy
equivalence. Thus the double complexes formed with the stalk complex compute
$\Ext_R^*(M, \mathcal{A}_x)$, and the maps from the neighbourhood complexes
give a comparison from the spectral sequence of Equation
(\ref{equation-constant-sheaf-ext-stalk-first-spectral-sequence}) to this
stalk spectral sequence. Its $E_2$-page is concentrated in the row $q = 0$
and that row is $\Ext_R^p(M, \mathcal{A}_x)$. Under hypothesis (1), the map
on $E_2$-pages is an isomorphism. Therefore it is an isomorphism on every
later page and on the graded pieces of the finite abutment filtrations.
Descending induction on those finite filtrations proves that the map on
abutments is an isomorphism. By construction that map is the canonical map
in Equation (\ref{equation-constant-sheaf-ext-stalk-comparison}).

Under hypothesis (2), use instead the spectral sequence of Equation
(\ref{equation-constant-sheaf-ext-stalk-quotient-spectral-sequence}). The
quotient hyper-Ext construction applied to the stalk complex has
$$
E_2^{p, q} =
\Ext_S^p(\text{Tor}_q^R(S, M), \mathcal{A}_x)
\quad\Longrightarrow\quad
\Ext_R^{p + q}(M, \mathcal{A}_x),
$$
because the stalk complex has cohomology $\mathcal{A}_x$ in degree zero and
zero in higher degrees. The displayed colimit calculation and hypothesis
(2) identify the two $E_2$-pages.
The same finite-filtration argument proves that the canonical map on
abutments is an isomorphism.
\end{proof}


\begin{lemma}
\label{lemma-constant-sheaf-ext-stalk-noetherian-quotient}
\begin{reference}
\cite[Chapter IV, \S 4.3, Remark 2(b), p. 193]{Tohoku}
\end{reference}
Let $X$ be a topological space, let $R$ be a ring, let $I \subset R$ be a
two-sided ideal, and put $S = R/I$. Let $M$ be an $R$-module, let
$\mathcal{A}$ be an $\underline{S}$-module on $X$, and let $x \in X$.
Assume that $S$ is left Noetherian and that
$\text{Tor}_a^R(S, M)$ is a finite $S$-module for every $a \geq 0$. Then
the canonical maps
$$
\SheafExt^n_{\underline{R}}(\underline{M}, \mathcal{A})_x
\longrightarrow \Ext_R^n(M, \mathcal{A}_x)
$$
are isomorphisms for all $n \geq 0$.
\end{lemma}

\begin{proof}
Put $T_a = \text{Tor}_a^R(S, M)$. Since $S$ is left Noetherian and $T_a$
is finite, $T_a$ has a resolution by finite free $S$-modules. Indeed, choose
a surjection from a finite free module onto $T_a$; its kernel is finite by
the left Noetherian hypothesis, and iteration gives the desired resolution.

Let $(N_i)$ be any filtered system of $S$-modules and let $P_\bullet$ be
such a resolution of $T_a$. Since every $P_j$ is finite free, the canonical
maps
$$
\colim_i \Hom_S(P_j, N_i)
\longrightarrow \Hom_S(P_j, \colim_i N_i)
$$
are isomorphisms. Filtered colimits of abelian groups are exact, so taking
cohomology of the Hom complexes gives isomorphisms
$$
\colim_i \Ext_S^p(T_a, N_i)
\longrightarrow \Ext_S^p(T_a, \colim_i N_i)
$$
for every $p \geq 0$. Apply this with
$N_U = H^b(U, \mathcal{A})$, where $\mathcal{N}_x$ is the filtered set of
open neighbourhoods of $x$ ordered by reverse inclusion. Hypothesis (2) of
Lemma \ref{lemma-constant-sheaf-ext-stalk-comparison} holds for every
$a, b, p$, and the result follows from that lemma.
\end{proof}
